\documentclass[12pt,a4paper]{article}
\usepackage[utf8]{inputenc}
\usepackage[french]{babel}
\usepackage{listings} % Permet d'inclure ton code C automatiquement
\usepackage{xcolor}   % Pour colorer le code

% Configuration des couleurs pour le code C
\definecolor{commentColor}{rgb}{0,0.6,0}
\definecolor{keywordColor}{rgb}{0,0,1}
\definecolor{stringColor}{rgb}{0.58,0,0.82}

\lstset{
    commentstyle=\color{commentColor},
    keywordstyle=\color{keywordColor},
    stringstyle=\color{stringColor},
    basicstyle=\footnotesize\ttfamily,
    breaklines=true,
    numbers=left,
    numbersep=5pt,
    language=C
}

\begin{document}

\title{Projet ASD : Structures de Données Hybrides}
\author{Seydina Mohamed Diop \& Abdou Lahate}
\date{\today}
\maketitle

\section{Introduction}
Ce Projet porte sur un systeme de gestion de scolarité .
En efet on va consevoire une programme qui va permetre a un utulisateur de rensegne
quelque information qui saffichera sur le terminal 

\section{Modélisation des Structures de Données}
Voici les Structures de base que nous avons conçue pour notre \texttt{Systeme De Scolarité}.

% 1. On affiche uniquement la structure Note (lignes 4 à 10 par exemple)
\lstinputlisting[linerange=4-10]{../structure.h}

% 2. Tu écris ton explication de la structure Note juste ici
\subsubsection{Explication de la structure Note}
La structure \texttt{Note} reçoit et stocke les resultats academiques d'un etudiant
pour chaque matière de la filiere, ainsi que le coefficient correspondant. 
Par exemple, pour un étudiant donné, cette structure va enregistrer ses notes au fur et à mesure. 
Ensuite, grace au pointeur \texttt{struct Note *note\_suivante}, la structure se connecte à la note
suivante du même étudiant.C'est ce mécanisme qui forme une liste chaînée de notes propre à chaque etudiant, 
permettant de stocker un nombre dynamique de modules sans utiliser de tableau statique.

\lstinputlisting[linerange=12-20]{../structure.h}

\subsubsection{Explication de la structure Etudiant}
La structure \texttt{Etudiant} regroupe l'ensemble des informations d'un étudiant de la filière, telles que son nom, son 
prénom, son numéro de dossier (identifiant unique) et son état d'inscription (via un booleen).
Le pointeur \texttt{modules}  est essentielle : il pointe vers la première note de l'etudiant, liant ainsi l'etudiant 
à ses résultats.Concernant le pointeur \texttt{etudiant\_suivant} elle pointe sur l etudiant suivante creent ainsi 
une liste chaine

\lstinputlisting[linerange=21-27]{../structure.h}

\subsubsection{Explication de la structure Filiere}
Cette structure contient le nom d'une filiere donner ainsi que le nombre d'etudiants.
Dans la structure \texttt{Filiere} on a deux pointeurs : \texttt{premier\_du\_liste} qui pointe sur le premier étudiant 
de la liste du filière, \texttt{dernier\_du\_liste} quant a lui pointe sur le dernier étudiant. Permettant d inserer des
etudiants directemant a la fin de la liste 

\section{Modélisation des Fonctions}
Voici les fonctions que nous avons conçue pour permetre le bon fonctionement de notre Systeme de scolarité.

\lstinputlisting[linerange=8-44]{../fonctions.c}

\subsubsection{Explication de la Fonctions inscrire\_etudiant}
Cette fonction nous permet d ajouter de nouveaux etudiant dans une filiere en enregitrant leur information personnel: prenon, non,numero de dossier  
grace a l allocation dynamique suivant:\texttt{Etudiant *nouveau\_etudiant = (Etudiant*) malloc(sizeof(Etudiant))} qui allout une espace memoire dans le memoire RAM.
Combiner au boucle suivant.
\lstinputlisting[linerange=30-39]{../fonctions.c}
l utulisateur poura ajouter autant d etudiant qu il souhaite, le programme aura ainsi le nombre totale d etudiants.



\lstinputlisting[linerange=45-111]{../fonctions.c}

\subsubsection{Explication de la Fonctions note\_etudiant}
Comme la fonction precedente celui ci Permet de stocke les resultats scolaire de chaque etudiants enregistrer en amont ,grace a l allocation dynamique suivant
\texttt{ Note *nouveau\_note = (Note*)malloc(sizeof(Note))}.CEtte fonction va permetre a l utilisateur de donner autant de matiere ainsi que les note et coefficient associer.
Pour chacune des etudiants qu il enregistre il peux aussi ne pas ajouter de note pour les etudiants qu il souhaite . Et cela grace au boucle \texttt{ while (nouveaux\_matiere == 2)}
qui lui dira de taper 2 s il veux ajouter des matiere et 0 sinon .





\lstinputlisting[linerange=113-128]{../fonctions.c}

\subsubsection{Explication de la Fonctions liste\_de\_la\_classe}
Quant a cette fonction il affichera la liste des etudiants que l utulisateur a deja enregistrer .Il donne respectivement leur nom, prenom et numero de dossier.






\lstinputlisting[linerange=130-191]{../fonctions.c}

\subsubsection{Explication de la Fonctions bulletin}
La fonction \texttt{bulletin} permet de calculer les moyenne de chaque etudiant ainsique la moyenne gennerale de la classe.
Elle va aller chercher les module de chaque etudiants c est a dire ces matiere ainsi que les coef et notes associer pour calculer sa moyenne.
Si l etudiant n a pas de note la boucle :
\lstinputlisting[linerange=171-175]{../fonctions.c}
Va afficher a l ecrant que cette etudiant n a pas de note ,cette etudiant ne sera pas prise en compte pour le calcule du moyenne gennerale de la classe .
Grace au variable \texttt{etudiants\_avec\_notes} , innicialise a 0 qui n avance que si l etudiant qu on doit calculer la moyenne a une note .De cette facon avec 
la boucle :
\lstinputlisting[linerange=182-191]{../fonctions.c}
Cqlcule et affiche la moyenne gennerale exacte de la classe.






\lstinputlisting[linerange=193-274]{../fonctions.c}

\subsubsection{Explication de la Fonctions recherche\_et\_modification\_etudiant}
Cette fonction va demander a l utilisateur s il veut rechercher un etudiant affin de modifier 
ou non ces information personnel ou academiques en l affichant un message lui permettantde faire deux choix
soit taper 3 pour contunuer et avoire la possibiliter de modifiers certainne information .Soit taper 0 pour arreter le programme.
et eventuellement ne rien modifier ou sortire apres des modification.







\lstinputlisting[linerange=276-335]{../fonctions.c}

\subsubsection{Explication de la Fonctions suprimer\_etudiant}
De la meme manier que la fonction precedente celui ci lui permetra de suprimer ou non un etudiant de la liste 








\lstinputlisting[linerange=337-363]{../fonctions.c}

\subsubsection{Explication de la Fonctions maximum}
Cette fonction va chercher parmit tout les etudiant celui qui a la plus grande moyenne et a  la fin du programme il 
donne sont nom prenom numero de dossier sa moyenne et le nom de la filiere.








\lstinputlisting[linerange=365-391]{../fonctions.c}

\subsubsection{Explication de la Fonctions minimum}
Pareil pour la fonction \texttt{minimum} ,il donne l etudiant ayant la plus faible moyenne de la filiere .










\lstinputlisting[linerange=393-447]{../fonctions.c}

\subsubsection{Explication de la Fonctions du\_majorant\_minorant}
Cette fonction Permet d afficher par ordre de merite tout les etudiant de la filiere du majorant a l etudiant ayant la plus faible moyenne
grace a la \texttt{trie par selection}:
\lstinputlisting[linerange=410-446]{../fonctions.c}









\lstinputlisting[linerange=449-501]{../fonctions.c}

\subsubsection{Explication de la Fonctions orde\_alphabet}
Consernant cette fonction elle utilise une \texttt{trie a bulle}, et elle va permetre l affichage de la liste
par ordre alphabetique :
\lstinputlisting[linerange=467-500]{../fonctions.c}







\section{Modélisation du corps de notre projet : (\texttt{int main})}
C est ici que nous faisont les appelles des differentes fonctions pour donner vie a notre programme.

\lstinputlisting[linerange=6-45]{../main.c}











\end{document}